\documentclass[10pt,a4paper]{amsart}
\usepackage[margin=19mm]{geometry}
\usepackage{amsmath,amssymb,mathtools}
\usepackage[scr=rsfs,cal=euler]{mathalfa}
\usepackage{xfrac}
\usepackage[unicode,hidelinks,bookmarks=false]{hyperref}
\newcommand{\ie}{i.e.\ }
\newcommand{\cf}{cf.\ }
\newcommand{\resp}{respectively\ }
\newcommand{\Subsection}{\subsection}
\providecommand{\nobreakdash}{\nobreak\hbox{-}}
\setlength{\emergencystretch}{2em}
\multlinegap0pt
\usepackage{dsfont}
\usepackage[shortlabels]{enumitem}
\usepackage{bm}
\usepackage{mathtools}
\usepackage{amssymb}
\newtheorem{lemma}{Lemma}
\newcommand{\Na}{\mathbb{N}}
\newcommand{\op}{\bm}
\newcommand{\spec}{\textttup{spec}}
\newcommand{\textttup}[1]{\textup{\texttt{#1}}}
\title{Operator-Theoretic Foundations and Policy Gradient Methods for General MDPs with Unbounded Costs}
\author{Abhishek Gupta, Aditya Mahajan}
\date{Source: arXiv:2603.17875v3 (CC BY 4.0)}
\begin{document}
\maketitle

\noindent\emph{Source and licence.} Operator-Theoretic Foundations and Policy Gradient Methods for General MDPs with Unbounded Costs. Version arXiv:2603.17875v3 (CC BY 4.0), distributed by arXiv under the Creative Commons Attribution 4.0 International licence, http://creativecommons.org/licenses/by/4.0/, as recorded in the arXiv licence metadata for this exact version and shown on the abstract page https://arxiv.org/abs/2603.17875v3. Full licence notice: \texttt{licenses/arxiv-2603.17875-CC-BY-4.0.txt}.\par\smallskip

\noindent\textit{Editorial scope.} Counted unit: the article's lemma lem:specA (source0.tex lines 1049--1051). The article's complete original proof of it is delivered with it (source0.tex lines 1052--1065, one \texttt{proof} environment of 14 lines). No mathematics is written, repaired or re-derived: the statement and the proof are the article's own lines, in the article's own notation, and no retained line is substituted or re-flowed.  Retained uncounted as the source-local context the proof needs: lines 288--288.  Nothing was omitted from any retained span, so the package carries no omission record.  No retained line contains a citation, so the wrapper carries no bibliography and no bibliography entry.  No retained line contains a figure environment, an image-inclusion command or drawing code, so no image is delivered and the package declares no asset dependency.  Editorial formatting only, in addition: an AMS \texttt{amsart} preamble that restates the article's own declarations and macros used by the retained lines, namely the environments \texttt{lemma} on separate counters, together with the standard packages those lines need.  The retained environments print in this document as Lemma 1 in order of appearance, whereas the article prints the same items with the numbering of its own full document.  The complete native source of the article as distributed in the arXiv e-print for this exact version is bundled unchanged as \texttt{sources/196655/source0.tex}.
\subsection{Primer on Linear Operators}\label{sub:linearop}

\begin{lemma}\label{lem:specA}
     If $\spec(\op A)<1$, then there exists constants $\tilde\kappa, \kappa>0$ such that for any $x\in\mathcal X$ satisfying $ \|x\|\leq \tilde\kappa\|y\|$, we have $\|\op B^k x\|\leq \kappa \|y\|$ for all $k\in \Na$.
\end{lemma}

\begin{proof}
If $\spec(\op A)<1$, then there exists $L\in\Na$ such that $\alpha:=\|\op A^L\|<1$ (see Subsection \ref{sub:linearop}).
Let $\beta = \sum_{i=0}^{L-1} \|\op A^i\|\in(0,\infty)$, $\tilde \kappa = \frac{\beta}{1-\alpha}>0$, and $\kappa = \beta+\tilde\kappa \max\{1,\|\op A\|,\ldots,\|\op A^{L-1}\|\}$. We now apply the principle of mathematical induction to complete the proof. First, we have an intermediate result: $\|\op B^{nL} x\| \leq \tilde \kappa \|y\|$ for all $n\in\Na$, which we establish through induction. Indeed, for $n=1$, if $ \|x\|\leq \tilde\kappa\|y\|$, then we have 
\[\|\op B^L x\|\leq \left\|\sum_{i=0}^{L-1} \op A^i y\right\| + \|\op A^L\| \|x\| \leq \beta \|y\| + \alpha \|x\|\leq \|\leq\beta \|y\| + \alpha \tilde \kappa \|y\| = \tilde \kappa \|y\|.\]
Suppose that this statement holds for $n\in\Na$. We use the same argument as above to get 
\begin{align*}
    \|\op B^{(n+1)L} x\| = \|\op B^L \op B^{nL}x\|\leq \left\|\sum_{i=0}^{L-1} \op A^i y\right\| + \|\op A^L\| \|\op  B^{nL} x\| \leq \tilde \kappa \|y\|. 
\end{align*}
The intermediate result in thus established. Let us establish the result now. For $1\leq l\leq L-1$, since $ \|x\|\leq \tilde \kappa\|y\|$, we have 
\begin{align}
    \|\op B^l x\| \leq \beta\|y\| + \|\op A^l\| \|x\|\leq (\beta+\|\op A^l\|\tilde\kappa)\|y\|\leq \kappa\|y\|.\label{eqn:Blx}
\end{align}
Next, we observe that $\op B^{nL+l} x = \op B^l \op B^{nL} x$, and we apply \eqref{eqn:Blx} with $x$ replaced with $\op B^{nL} x$ to conclude that $\|\op B^{nL+l} x\|\leq \kappa \|y\|$. This completes the induction step and we arrive at the conclusion that $\|\op B^k x\|\leq \kappa \|y\|$ for all $k\in \Na$ (recall that by construction, $\kappa>\tilde\kappa$). 
\end{proof}

\end{document}
